% easy/easy.tex
% mainfile: ../perfbook.tex
% SPDX-License-Identifier: CC-BY-SA-3.0

\QuickQuizChapter{chp:Ease of Use}{易用性}{qqzeasy}
%
\Epigraph{创建一个完美的API就像实施一场完美的犯罪。
	  至少有五十件事情可能出错，即使你是天才，
	  也最多只能预见其中二十五件。}
	 {向所有可能还健在的Kathleen Turner粉丝致歉。}

\section{什么是"容易"？}
\label{sec:easy:What is Easy?}
%
\epigraph{当有人说"我想要一种编程语言，只需说出我想做什么就行"时，
	  给他一根棒棒糖。}
	 {Alan J.~Perlis，经改编}
% http://www.cs.yale.edu/homes/perlis-alan/quotes.html

如果你有轻视易用性需求的倾向，请考虑这样一个事实：
Linux内核RCU中的一个易用性bug导致了一个可被利用的
Linux内核安全漏洞~\cite{McKenney:2019:CRS:3319647.3325836}。
因此，即使是内核内部的API也必须易于使用，这一点显然很重要。

不幸的是，"容易"是一个相对的概念。
例如，许多人会认为15小时的飞机旅程有点折磨人——
除非他们停下来考虑替代的交通方式，尤其是游泳。
这意味着创建易用的API要求你足够了解目标用户，
知道什么对他们来说是容易的。
而这可能与什么对你来说是容易的毫无关系。

以下问题说明了这一点：
"从当今世界上随机选择一个人，做出什么样的单一改变能改善此人的生活？"

没有任何单一的改变能保证帮助每个人的生活。
毕竟，人与人之间差异极大，需求、愿望、渴望和抱负也相应地千差万别。
一个饥饿的人可能需要食物，但额外的食物可能会加速一个病态肥胖者的死亡。
许多年轻人如此热切渴望的高度刺激，对于心脏病恢复期的人来说可能是致命的。
对某个人的成功至关重要的信息，可能会加剧另一个遭受信息过载之苦的人的失败。
简而言之，如果你正在开发一个旨在帮助你一无所知的人的软件项目，
当这些人对你的项目提出批评时，你不应该感到惊讶。

如果你真的想帮助特定的一群人，
没有什么能替代与他们长期紧密合作，以年为计。
尽管如此，有一些简单的事情可以增加用户对你的软件满意的可能性，
其中一些将在下一节中介绍。

\section{API设计的Rusty量表}
\label{sec:easy:Rusty Scale for API Design}
%
\epigraph{因此，找到合适的度量不是一个数学练习，
	  而是一个冒险的判断。}
	 {Peter Drucker}
% http://billhennessy.com/simple-strategies/2015/09/09/i-wish-drucker-never-said-it
% Rusty is OK with this: July 19, 2006.

本节改编自Rusty Russell 2003年渥太华Linux研讨会主题演讲的
部分内容~\cite[Slides~39--57]{RustyRussell2003OLSkeynote}。
Rusty的关键观点是，目标不应该仅仅是使API易于使用，
而应该是使API难以被误用。
为此，Rusty提出了他的"Rusty量表"，
按照这一重要的"难以误用"属性的递减顺序排列。

以下列表尝试将Rusty量表推广到Linux内核之外：

\begin{enumerate}
\item	不可能用错。
	虽然这是所有API设计者应该努力达到的标准，
	但只有神话般的\co{dwim()}\footnote{
		\co{dwim()}函数是一个首字母缩写，
		展开为"do what I mean"（做我想做的）。}
	命令才能接近这个目标。
\item	编译器或链接器不会让你用错。
\item	如果你用错了，编译器或链接器会发出警告。
	\co{BUILD_BUG_ON()}是用户的好朋友。
\item	最简单的用法就是正确的用法。
\item	名称告诉你如何使用它。
	但名称可能是双刃剑。
	虽然\co{rcu_read_lock()}对于从reader-writer locking
	转换代码的人来说足够明了，
	但对于从引用计数转换代码的人来说可能会引起一些困惑。
\item	用对了就没事，用错了运行时一定会出问题。
	\co{WARN_ON_ONCE()}是用户的好朋友。
\item	遵循常见惯例，你就能用对。
	\co{malloc()}库函数就是一个好例子。
	虽然很容易在内存分配上出错，
	但大量项目确实在大多数时候做到了正确使用。
	将\co{malloc()}与Valgrind~\cite{ValgrindHomePage}
	配合使用，几乎可以将\co{malloc()}提升到
	量表上"用对了就没事，用错了运行时一定会出问题"的位置。
\item	阅读文档，你就能用对。
\item	阅读实现代码，你就能用对。
\item	阅读正确的邮件列表归档，你就能用对。
\item	阅读正确的邮件列表归档，你还是会用错。
\item	阅读实现代码，你还是会用错。
	\co{rcu_read_lock()}的原始非\co{CONFIG_PREEMPT}
	实现~\cite{PaulEMcKenney2007PreemptibleRCU}
	就是量表上这一位置的臭名昭著的例子。
\item	阅读文档，你还是会用错。
	例如，DEC Alpha的\co{wmb}指令的
	文档~\cite{ALPHA2002}让许多开发者
	误以为该指令具有比实际强得多的内存顺序语义。
	后来的文档澄清了这一点~\cite{Compaq01,WilliamPugh2000Gharachorloo}，
	将\co{wmb}指令提升到了量表上
	"阅读文档，你就能用对"的位置。
\item	遵循常见惯例，你还是会用错。
	\co{printf()}语句就是量表上这一位置的例子，
	因为开发者几乎总是忘记检查\co{printf()}的错误返回值。
\item	用对了，运行时反而会出问题。
\item	名称告诉你的是如何\emph{不}使用它。
\item	显而易见的用法是错误的。
	Linux内核的\co{smp_mb()}函数就是量表上这一位置的例子。
	许多开发者认为这个函数具有比其实际拥有的强得多的
	顺序语义。
	\Cref{chp:Advanced Synchronization: Memory Ordering}包含了
	避免这一错误所需的信息，Linux内核源码树的
	\path{Documentation}和\path{tools/memory-model}目录也是如此。
\item	如果你用对了，编译器或链接器会发出警告。
\item	编译器或链接器不会让你用对。
\item	不可能用对。
	\co{gets()}函数就是量表上这一位置的著名例子。
	事实上，\co{gets()}最好的描述可能是
	一个无条件的缓冲区溢出安全漏洞。
\end{enumerate}

\section{修剪曼德博集合}
\label{sec:easy:Shaving the Mandelbrot Set}
%
\epigraph{简单不在复杂之前，\\而在复杂之后。}
	 {Alan J.~Perlis}

有用程序的集合类似于曼德博集合
（如\cref{fig:easy:Mandelbrot Set}所示），
因为它没有清晰光滑的边界——如果有的话，
停机问题就可以求解了。
但我们需要的是真实的人能够使用的API，
而不是每一个潜在用途都需要完成一篇博士论文的API。
因此，我们"修剪曼德博集合"\footnote{
	归功于Josh Triplett。}，
将API的使用限制在完整潜在用途集合中一个容易描述的子集内。

\begin{figure}
\centering
\resizebox{2.5in}{!}{\includegraphics{easy/Mandel_zoom_00_mandelbrot_set}}
\caption{曼德博集合（图片来自维基百科）}
\label{fig:easy:Mandelbrot Set}
\end{figure}

这样的修剪可能看起来适得其反。
毕竟，如果一个算法能工作，为什么不应该使用它呢？

要理解为什么至少某些修剪是绝对必要的，
考虑一个避免\IX{deadlock}的locking设计，但可能以最糟糕的方式。
这个设计使用一个循环双向链表，
其中包含系统中每个thread的一个元素以及一个头元素。
当一个新thread被创建时，父thread必须向这个列表中插入一个新元素，
这需要某种形式的同步。

保护列表的一种方法是使用全局lock。
然而，如果thread被频繁创建和删除，
这可能会成为瓶颈。\footnote{
	有操作系统背景的读者，请暂时搁置疑虑。
	无法搁置疑虑的读者欢迎提供更好的例子。}
另一种方法是使用hash table并锁定单独的hash桶，
但当按顺序扫描列表时，这可能表现不佳。

第三种方法是锁定单独的列表元素，
并要求在插入时同时持有前驱和后继元素的lock。
由于必须获取两把lock，我们需要决定获取的顺序。
两种传统方法是按地址顺序获取lock，
或按它们在列表中出现的顺序获取lock，
这样当头元素是被锁定的两个元素之一时，它总是先被获取。
然而，这两种方法都需要特殊的检查和分支。

这个待修剪的方案是无条件地按列表顺序获取lock。
但deadlock怎么办？

Deadlock不可能发生。

要理解这一点，将列表中的元素从零开始编号，
头元素为零，直到列表中的最后一个元素$N$
（即在循环列表中位于头元素之前的那个）。
类似地，将thread从零编号到$N-1$。
如果每个thread试图锁定某对连续的元素，
至少有一个thread能够保证获取两把lock。

为什么？

因为没有足够多的thread来环绕整个列表。
假设thread~0获取了元素~0的lock。
要使其被阻塞，某个其他thread必须已经获取了元素~1的lock，
所以让我们假设thread~1已经这样做了。
类似地，要使thread~1被阻塞，某个其他thread必须已经获取了
元素~2的lock，依此类推，直到thread~$N-1$获取了
元素~$N-1$的lock。
要使thread~$N-1$被阻塞，某个其他thread必须已经获取了
元素~$N$的lock。
但已经没有更多的thread了，所以thread~$N-1$不可能被阻塞。
因此，deadlock不可能发生。

那么为什么我们应该禁止使用这个妙趣横生的小算法呢？

事实是，如果你\emph{真的}想使用它，我们无法阻止你。
然而，我们\emph{可以}建议不要将此类代码包含在
我们关心的任何项目中。

但是，在你使用这个算法之前，
请仔细思考以下Quick Quiz。

\QuickQuiz{
	在删除元素时可以使用类似的算法吗？
}\QuickQuizAnswer{
	可以。
	然而，由于每个thread必须持有三个连续元素的lock
	才能删除中间的那个，如果有$N$个thread，
	则必须有$2N+1$个元素（而不仅仅是$N+1$个）
	才能避免deadlock。
}\QuickQuizEnd

事实是，这个算法极其专门化（它只对特定大小的列表有效），
而且相当脆弱。
任何意外未能将节点添加到列表的bug都可能导致deadlock。
实际上，仅仅是稍微晚了一点添加节点就可能导致deadlock，
增加thread的数量也是如此。

此外，上面描述的其他算法是"好且足够的"。
例如，简单地按地址顺序获取lock相当简单且快速，
同时允许使用任意大小的列表。
只要注意空列表和只包含一个元素的列表所带来的特殊情况！

\QuickQuiz{
	呸！
	到底是什么人想出了一个如此值得被修剪的算法？？？
}\QuickQuizAnswer{
	那就是Paul。

	他当时在思考\emph{哲学家就餐问题}，
	这涉及一顿由五位哲学家参加的相当不卫生的意大利面晚餐。
	考虑到桌上有五个盘子但只有五把叉子，
	而且每位哲学家每次吃饭需要两把叉子，
	要求想出一种避免deadlock的叉子分配算法。
	Paul的回应是"天哪！
			  再拿五把叉子不就行了！"

	这本身没什么问题，但Paul随后将同样的方案
	应用到了循环链表上。

	这也不算太糟糕，但他偏偏还要告诉别人！
}\QuickQuizEnd

\begin{figure}
\centering
\resizebox{2.5in}{!}{\includegraphics{cartoons/r-2014-shaving-the-mandelbrot}}
\caption{修剪曼德博集合}
\ContributedBy{fig:easy:Shaving the Mandelbrot Set}{Melissa Broussard}
\end{figure}

总而言之，我们不会仅仅因为算法碰巧能工作就使用它。
相反，我们将自己限制在那些足够有用、
值得花时间学习的算法上。
算法越困难、越复杂，它就必须越具有通用性，
才能使学习它和修复其bug的痛苦变得值得。

\QuickQuiz{
	给出这条规则的一个例外。
}\QuickQuizAnswer{
	一个例外是，一个困难而复杂的算法，
	却是在给定情况下已知的唯一可行算法。
	另一个例外是，一个困难而复杂的算法，
	却是在给定情况下已知可行算法集合中最简单的一个。
	然而，即使在这些情况下，
	花一点时间尝试想出一个更简单的算法也可能非常值得！
	毕竟，如果你能够发明出第一个完成某项任务的算法，
	那么继续发明一个更简单的算法应该也不会太难。
}\QuickQuizEnd

除了例外情况，我们必须继续修剪软件的"曼德博集合"，
使我们的程序保持可维护性，
如\cref{fig:easy:Shaving the Mandelbrot Set}所示。

\QuickQuizAnswersChp{qqzeasy}
